\section{Stopped local sets and canonical clusters}
\label{sec:local-sets}

We now identify the field left by a stopped current exploration. The
discrete field in this section is \(h_\delta=aK_\delta\), whose limit has
the normalization of \(h\). Two kinds of stops are needed. Stopping at a
fixed number of odd rims gives the two-valued carpets. A seed exploration
instead reveals whole current clusters; its limiting conditional field
will be used to decide which odd rims belong to one cluster. For both
stops we prove the conditional law on tests that cross the limiting cut.
This is stronger than identifying the field only in its open components.

A closed set \(A\), together with information \(\mathscr E\), is a local
set of \(h\) when, conditionally on \(\mathscr E\),
\(h=h_A+h^{D\setminus A}\), where \(h_A\) is
\(\mathscr E\)-measurable and harmonic off \(A\), and
\(h^{D\setminus A}\) is the zero-extended GFF with independent
zero-boundary restrictions to the components of \(D\setminus A\).
Conditioning down to \(\sigma(A,h_A)\) gives the usual local-set definition.
It is \emph{thin} when \(h_A\) is integration against that harmonic
function, with no further distribution supported on \(A\). We include
the domain boundary in every closed local set. Thus statements about
zero area below concern its intersection with the open domain; a Jordan
boundary itself need not have zero area.

For \(p,q>0\) with \(p+q\ge2\lambda\), a boundary-connected thin local
set with complementary values in \(\{-p,q\}\) is the canonical
\(\TVS{p}{q}\), by \cite[Proposition~2]{ASW2017}. Here
boundary-connected means that the closed set, with the boundary adjoined,
is connected; it supplies the topological condition in the bounded-type
thin-local-set theorem. This uniqueness does not assume that the candidate
was already measurable in the field. The resulting canonical set is
field-measurable. We use the Jordan-hole property in Jordan parents. When
a first odd carpet is constructed in a rough simply connected parent, its
two label magnitudes are \(2a>2\lambda\); its loops avoid the parent
boundary and hence remain Jordan under the interior conformal map.
These geometric facts and the nested constructions are in
\cite[Sections~3--4]{AS2018}. In particular
\(\TVS{2\lambda}{2\lambda}\) has \(\CLE_4\) holes with independent fair
signs and conditionally independent zero-boundary remainders.

\subsection{Stops and their cut neighborhoods}

We first specify the finite stops. In a free-current chamber of flat value
\(b\), descend through surrounding current clusters, continue through even
holes, and stop on the first odd rim along each branch. The closed complement
of the unsearched holes is the \emph{first odd carpet}; its hole values are
\(b-2a\) and \(b+2a\). Repeating this operation in the holes defines the
odd depth. This counts odd rims, not all clusters inspected: any number of
even holes may be passed before the next odd rim. All holes are the sealed
component faces of \cref{lem:finite-cuts}, with its medial resolution at
pinches.

For a seed stop, use a fixed predictable growth rule: from the current
revealed data choose the next vertex and an attachment to the revealed
region, reveal the entire current cluster first met there, and stop using
the revealed data. Apart from the new vertex, the attachment uses only
already revealed vertices or known seed edges, with no unsearched free
vertex in between. The rule may follow a preassigned column sweep or choose
the lexicographically first unvisited vertex adjacent to the revealed
region. It may use discovery indices, observed diameters, or the first
observed odd jump around specified points. The seed trace may therefore
be adaptive; its union with the revealed clusters remains connected to
the exterior.
The terminal sigma-field \(\mathcal E_\delta\) includes the revealed
spins, edges, colors, heights, and the seed. The initial sigma-field
\(\mathcal E_{0,\delta}\) records only the exact flat cut from which this
exploration starts. Conditional on it, the initial unsearched chambers
have the independent flat laws of \cref{lem:finite-cuts}.

We also use a composite seed stop for the wired current. Expose its
boundary cluster first, and then run one fixed global seed rule through
the remaining free chambers, attaching each newly visited seed piece
through the already revealed region containing that layer. A half-column sweep can be ordered
this way. Seed edges are edges of the underlying graph and need not be
current-open. After every added vertex, reveal the whole cluster met
there. The revealed union remains connected to the original boundary,
and the finite cut law applies simultaneously to all its residual
chambers. This is one exploration of the original field, even when one
discrete wired hole will later separate into several limiting holes.

Write \(A_\delta\) for the revealed closed set with its one-cell
thickening, \(S_\delta\) for the seed at the stop, and \(W_{0,\delta}\)
for the initial revealed wall. For a carpet the seed is empty. Let
\(U_{0,\delta}\) be the union of the initial open chamber faces and let
\(F_{0,\delta}\) be its closed complement in the fixed enclosing disk
\(\overline B_0\) of \cref{sec:gaussian}. Define \(U_\delta,F_\delta\)
in the same way for the residual faces at the terminal stop. On compact
subsets of the initial domain, \(F_\delta\) and \(A_\delta\) differ only
by the \(O(\delta)\) medial displacement described immediately before
\cref{prop:conditional-flat-gff}.

Use the same atomic smearing for all three finite distributions
\[
 M_\delta(\phi)=\E[h_\delta(\phi)\mid\mathcal E_\delta],
 \qquad R_\delta=h_\delta-M_\delta,\qquad h_\delta=M_\delta+R_\delta.
\]
Thus \(R_\delta\) is \(a\) times the centered flat remainders of
\cref{prop:conditional-flat-gff}, patched by zero on exposed sites.
For bounded nonsmooth tests, a set denotes the union of its lattice
cells and the pairing is the corresponding atomic sum. This convention
also identifies it with piecewise constant integration up to the
smooth-test cell error.

For a discovered cluster \(C\), let \(\mathcal O_\delta(C)\) be the union
of its odd holes. A contributor to a compact set \(Q\) is a discovered
\(C\) for which \(C\) or \(\mathcal O_\delta(C)\) meets \(Q\); write
\(N_\delta(Q)\) for their number. On the event
\[
 d_\delta(Q):=\dist(Q,F_{0,\delta}\cup S_\delta)\ge\rho>0,
\]
each contributor either reaches \(Q\) from the seed or surrounds a point
of \(Q\), and has diameter at least a fixed multiple of \(\rho\).
Consequently \(N_\delta(Q)\) is pathwise bounded by the initial
whole-box count of free clusters of that diameter. The domination and
tail in \cref{cor:global-cluster-counts}, conditional only on
\(\mathcal E_{0,\delta}\), make these counts uniformly tight over the
stopping rules. A completed terminal history can already record many
contributors, so no count tail conditional on that history is used.

\begin{lemma}[Negligible cut neighborhoods]
\label{lem:cut-neighborhoods}
Let \(Q\) be a compact subset of the initial domain and fix \(\rho>0\).
Consider a seed stop just described, a carpet of fixed odd depth, or the
first wired layer. For the last two cases omit \(S_\delta\) in
\(d_\delta(Q)\). There are \(C,\alpha>0\), depending on \(Q,\rho\) and
the fixed depth when applicable, such that for \(0<\eta<\rho/20\),
\begin{equation}
 \E\!\left[
   \ind_{\{d_\delta(Q)\ge\rho\}}\,|Q\cap A_\delta^\eta|
       \,\middle|\,\mathcal E_{0,\delta}\right]
 \le C(\eta+\delta)^\alpha .
 \label{eq:cut-area}
\end{equation}
The bound is averaged over the production of the stop and is uniform
over the allowed initial flat data and stopping rules.

Let \(\zeta_\rho\) be a continuous function on \([0,\infty)\), with
values in \([0,1]\), equal to zero on \([0,\rho]\) and to one on
\([2\rho,\infty)\). For every bounded \(\phi\) supported in \(Q\),
\begin{equation}
 \lim_{\eta\downarrow0}\limsup_{\delta\downarrow0}
 \E\!\left[\zeta_\rho(d_\delta(Q))
       |R_\delta(\phi\ind_{A_\delta^\eta})|^2\right]=0.
 \label{eq:remainder-neighborhood}
\end{equation}
For a fixed-depth carpet the analogous first-moment statement holds for
\(M_\delta\). For a seed stop it holds after restricting to
\(N_\delta(Q)\le M\) and to a fixed bound on the initial flat labels
of the chambers meeting \(Q\):
\begin{equation}
 \lim_{\eta\downarrow0}\limsup_{\delta\downarrow0}
 \E\!\left[\zeta_\rho(d_\delta(Q))\ind_{\{N_\delta(Q)\le M\}}
       |M_\delta(\phi\ind_{A_\delta^\eta})|\right]=0.
 \label{eq:mean-neighborhood}
\end{equation}
The restriction on the initial labels is understood in this display.
For the initial free and wired constructions their magnitudes are
respectively \(0\) and \(a\). The contributor restriction can be
removed in probability using the initial count tail.
\end{lemma}

\begin{proof}
Fix the initial exact-cut data. Suppose a seed stop satisfies
\(d_\delta(Q)\ge\rho\) and its revealed set comes within \(\eta\) of
\(z\in Q\). Up to one lattice cell, a newly revealed cluster then meets
\(B(z,O(\eta+\delta))\) and contains a current path from there to its
seed vertex at distance at least \(\rho\). Thus this event is contained
in the event of an \(N\)-arm from the small disk to radius \(\rho\).
This inclusion holds for every terminal rule, including rules using
the transverse colors.

Test this arm in the \((s,W)\) marginal of the initial flat kernel.
In a stack of separated buffered annuli between these radii, every
prescribed-color \(W\) circuit blocks the arm. Reveal only \((s,W)\)
in successive annuli, with endpoints before edges. The conditional
bound of \cref{prop:circuits} gives probability at least \(p>0\) for
each blocking circuit. The arm probability is therefore at most
\((1-p)^{c\log(\rho/(\eta+\delta))}\). This calculation never
conditions that circuit estimate on the full terminal history.
Integrating the event inclusion over \(z\in Q\) proves
\eqref{eq:cut-area} for seed stops, including the product of initial
flat chambers in the composite construction.

For a carpet of odd depth \(r\), use successive pairs of annuli with
opposite prescribed \(W\) signs. Between the two circuits in each
successful pair, the jump rule of \cref{lem:finite-cuts} forces an odd
rim. Nested successful pairs force distinct odd layers. Reaching the
inner disk before depth \(r\) therefore implies fewer than \(r\)
successes among logarithmically many one-parity trials. Their
conditional lower bounds give a binomial tail; its polynomial factor
is absorbed by decreasing the exponent. For the first wired layer,
a \(W\) circuit blocks the boundary current's arm, using the bulk
pre-removal clause of \cref{prop:circuits}. These arguments take place
in the original or fresh exact-cut kernel before any future color
revelations.

Condition now on the terminal data, only to estimate the remainder.
The residual fields are independent flat fields by
\cref{lem:finite-cuts}. For every terminal-measurable union \(E\) of
cells in \(Q\), the logarithmic covariance bound gives
\begin{equation}
 \E[|R_\delta(\phi\ind_E)|^2\mid\mathcal E_\delta]
 \le C\|\phi\|_\infty^2
 \iint_{E\times E}\left(1+\log^+
          \frac{R_0}{|z-w|\vee\delta}\right)\dd z\,\dd w
 \le C'\|\phi\|_\infty^2|E|.
 \label{eq:random-neighborhood-green}
\end{equation}
The final bound follows from the uniform integral of the logarithmic
kernel on a bounded set; its atomic version includes the diagonal
cell cutoff. Multiplying by \(\zeta_\rho(d_\delta(Q))\), averaging, and
using \eqref{eq:cut-area} proves
\eqref{eq:remainder-neighborhood}.

The revealed values and wall labels at fixed odd depth are bounded
by a deterministic depth-dependent constant. At a seed stop, on the
stated truncation the exact jump rule instead gives
\[
        |M_\delta|\le B+2aM+a\quad\hbox{on }Q,
\]
where \(B\) bounds the initial labels. Each contributor gives either
zero or one jump of size \(2a\); on its own sites the extra parity
discrepancy is at most \(a\). The area estimate proves
\eqref{eq:mean-neighborhood}. The pathwise contributor bound above
and the initial whole-box tail remove its truncation in probability.
\end{proof}

\subsection{The conditional field across a stopped cut}

We specify what is retained when a finite stop has a limit. Let
\(\mathbf e_\delta\) be a random element of a countable product of
Polish spaces, measurable at that same terminal stop. It
contains \(A_\delta,F_\delta,F_{0,\delta},W_{0,\delta},S_\delta\) as
closed-set coordinates. For each rational closed ball \(B\) and each
positive rational \(c\), it also records the terminal chamber label
when \(\dist(B,F_\delta)>c\), and the initial chamber label when
\(\dist(B,F_{0,\delta})>c\), with a cemetery value for either coordinate
when its clearance fails. Rational polygonal paths with positive clearance encode
connection between such buffers. One may include specified terminal-measurable
contributor coordinates whose chosen encodings are tight, such as the
localized finite lists used in \cref{sec:identification}. This retained
state excludes future clusters and canonical data that may appear
elsewhere in a larger joint representation.

For a fixed \(s<-1\), pass to a joint subsequence in the product of
\(H^s_{\mathrm{loc}}(B_0)^3\) and this retained state space:
\begin{equation}
 (h_\delta,M_\delta,R_\delta,\mathbf e_\delta)
       \ \Longrightarrow\ (h,M,R,\mathbf e),
 \qquad h=M+R.
 \label{eq:stopped-joint-extraction}
\end{equation}
The field tightness and the zero-extension tightness of the remainders
give this extraction for the first three coordinates. For a fixed
protected ball choose a smooth unit-mass test supported inside it and
divide its atomic pairing by the deterministic atomic mass of that test.
The mass tends to one, and on the clearance event the normalized
\(M_\delta\) pairing is exactly the chamber label. The full pairing is
tight and the remainder pairing has uniformly bounded second moment, so
the label on that event, with the cemetery value off it, is tight as well.
The same argument at the initial cut uses its centered flat remainder
and proves tightness of the protected initial-label coordinates.
This makes a diagonal extraction possible for the countable state;
unprotected labels in chambers collapsing onto the cut are not retained.
Write
\(\mathscr E=\sigma(\mathbf e)\), completed under the limiting law.
This is the limiting stopped sigma-field used below.

Let \(A,F,F_0,W_0,S\) be the closed-set limits. The one-cell comparison
between \(A_\delta\) and \(F_\delta\) gives \(A\cap D=F\cap D\);
similarly \(W_0\cap D=F_0\cap D\). Put \(U=D\setminus A\).
Its components are defined geometrically from \(F\): two rational
points belong to the same one exactly when some rational polygonal
path between small neighborhoods of them has positive clearance.
They are not defined by a limiting discrete equivalence relation.
At a rational point outside \(F\), choose a small ball and a rational
threshold strictly below its limiting clearance. Its protected coordinate
reads the limiting label; the point receives the cemetery value
on \(F\). The labels agree within each component, since a compact
connecting path is eventually inside one discrete chamber. Denote this
label by \(m_O\) on a component \(O\).

\begin{lemma}[Full conditional field at a stopped cut]
\label{lem:stopped-full-field}
In \eqref{eq:stopped-joint-extraction}, let \(Q\Subset D\) be compact.
On the stopped-measurable event
\(\Lambda_Q=\{\dist(Q,F_0\cup S)>0\}\), the set \(A\cap Q\) has zero
area, \(M\) is the distribution given on \(Q\) by the locally
integrable function \(m_O\) on each component of \(U\), and, for every
real \(\phi\in C_c^\infty(D)\) supported in \(Q\),
\begin{equation}
 \ind_{\Lambda_Q}
 \E[\exp(i(h-M)(\phi))\mid\mathscr E]
 =\ind_{\Lambda_Q}
       \exp\!\left(-\tfrac12G_U(\phi,\phi)\right).
 \label{eq:stopped-full-characteristic}
\end{equation}
The identity holds for every real linear combination of any finite
family of such tests. Thus, on tests crossing \(A\) in this region,
\(h-M\) is conditionally the zero-extended GFF on \(U\), with zero
covariance between its components. The mean is the limit of the
finite jump contributions after the contributor and initial-label
truncations of \cref{lem:cut-neighborhoods}, followed by their removal.

For a carpet explored in a single initial parent, take \(D\) to denote
that parent; then \(\Lambda_Q\) holds for every \(Q\Subset D\), and hence
\(|A\cap D|=0\). For a global wired carpet whose initial wall already
contains \(W_0\), the conclusion is used on compact subsets of each
pointed wired hole and combined with the previously identified layer.
For seed and composite stops the conclusions hold locally away from
\(W_0\cup S\). Each stop in a countable family has
this conclusion conditional on its own retained state; adjoining the
other stops to a joint representation does not enlarge that conditioning
sigma-field.
\end{lemma}

\begin{proof}
Use a Skorokhod representation of the joint subsequence for this proof.
We first derive zero area before using any continuum neighborhood bound.
Let \(d_\delta=d_\delta(Q)\) and \(d=\dist(Q,F_0\cup S)\).
Their convergence follows from Hausdorff convergence. For fixed \(\eta\),
every point of \(A\cap Q\) eventually lies in \(A_\delta^\eta\).
Fatou's lemma and \eqref{eq:cut-area} therefore give
\[
 \E[\zeta_\rho(d)|A\cap Q|]
 \le \liminf_{\delta\downarrow0}
       \E[\zeta_\rho(d_\delta)|Q\cap A_\delta^\eta|]
 \le C\eta^\alpha.
\]
Letting \(\eta\downarrow0\) proves zero area where \(d\ge2\rho\).
Exhausting positive rational \(\rho\) proves it on \(\Lambda_Q\).
The same holds for \(F\cap Q\), because \(A\) and \(F\) agree in \(D\).

We next prove the conditional characteristic function across the cut.
Fix a small \(\eta>0\). Choose a finite nonnegative smooth partition
\((\psi_i)_{i\le n_\eta}\) equal to one in sum on \(Q\), with each
\(\psi_i\) supported in a rational ball and with a fixed larger rational
closed ball \(B_i\Subset D\) of diameter at most \(4\eta\). Choose a continuous
\(\chi:[0,\infty)\to[0,1]\) equal to zero on \([0,1/4]\) and to one on
\([1,\infty)\), and put
\begin{equation}
 w_{i,\delta}=\chi\!\left(\frac{\dist(B_i,F_\delta)}{\eta}\right),
 \qquad
 \phi_\delta^\eta=\sum_i w_{i,\delta}\phi\psi_i.
 \label{eq:stopped-partition}
\end{equation}
Use the corresponding \(w_i,\phi^\eta\) with \(F\) in the limit.
These weights are continuous functions of the closed-set coordinate.
If \(w_{i,\delta}>0\), the whole buffer \(B_i\) has clearance at least
\(\eta/4\). If \(1-w_{i,\delta}>0\), every point in
\(\supp(\phi\psi_i)\) is within \(5\eta\) of \(F_\delta\). Consequently
\(\phi-\phi_\delta^\eta\) is bounded by \(\|\phi\|_\infty\) and
supported in \(Q\cap F_\delta^{5\eta}\).

Stratify the finitely many possible active index sets
\(\{i:w_{i,\delta}>0\}\). On each stratum,
\cref{eq:conditional-flat-characteristic}, multiplied by \(a\) in the
field normalization, applies to the fixed tests \(\phi\psi_i\).
Its fixed real coefficients may be replaced by the stopped-measurable
vector \((w_{i,\delta})\): approximate the compact coefficient cube
by a finite rational net. The conditional second-moment bound makes
the characteristic function uniformly Lipschitz in those coefficients,
and \(G_{U_\delta}\le G_{B_0}\) gives the same uniform continuity for
the Gaussian expression. Summing over the active strata yields
\begin{equation}
 \E\left|
 \E[e^{\,iR_\delta(\phi_\delta^\eta)}\mid\mathcal E_\delta]
       -e^{-\frac12G_{U_\delta}(\phi_\delta^\eta,\phi_\delta^\eta)}
       \right|\longrightarrow0.
 \label{eq:stopped-buffer-characteristic}
\end{equation}
This used only terminal-measurable coefficients and ancillaries, never
information from a later interior revelation.

For indices with positive limiting weight, their buffers retain positive
clearance. Green convergence in \cref{lem:flat-green-geometry} applies
to every pair of their fixed tests. Indices whose weights vanish are
harmless by the uniform bound from \(G_{B_0}\). Hence the Gaussian
exponent in \eqref{eq:stopped-buffer-characteristic} converges to
\(G_U(\phi^\eta,\phi^\eta)\). The left test converges jointly by
\eqref{eq:stopped-joint-extraction} and the finite continuous weights.
Multiplying by a bounded continuous function of \(\mathbf e_\delta\)
and by \(\zeta_\rho(d_\delta)\) is therefore legitimate.

There are two neighborhood errors to remove. The one-cell relation
between \(F_\delta\) and \(A_\delta\), \eqref{eq:random-neighborhood-green},
and \eqref{eq:cut-area} imply
\[
 \lim_{\eta\downarrow0}\limsup_{\delta\downarrow0}
 \E[\zeta_\rho(d_\delta)
       |R_\delta(\phi-\phi_\delta^\eta)|^2]=0.
\]
For the limiting kernel, positivity and domain monotonicity give
\begin{equation}
 \|\phi-\phi^\eta\|_{G_U}^{\,2}
 :=G_U(\phi-\phi^\eta,\phi-\phi^\eta)
 \le C\|\phi\|_\infty^2|Q\cap F^{5\eta}|.
 \label{eq:stopped-green-removal}
\end{equation}
Indeed \(0\le G_U\le G_{B_0}\), whose logarithmic singularity has a
uniformly bounded integral. The right side tends to zero on
\(\Lambda_Q\) by the zero-area result. Energy Cauchy--Schwarz now shows
that \(G_U(\phi^\eta,\phi^\eta)\to G_U(\phi,\phi)\) there; the other
energy factors are bounded by the same outer kernel.

The inequality \(|e^{ix}-e^{iy}|\le|x-y|\) removes the lattice error
in the tested conditional identity. Then
\eqref{eq:stopped-green-removal} and bounded convergence remove the
continuum error. We obtain, for every bounded continuous \(Z\) of the
retained state,
\[
 \E[\zeta_\rho(d)Z(\mathbf e)e^{\,iR(\phi)}]
   =\E[\zeta_\rho(d)Z(\mathbf e)e^{-\frac12G_U(\phi,\phi)}].
\]
Bounded continuous functions determine the law on the retained metric
state. A monotone-class argument gives the same equality for all bounded
\(\mathscr E\)-measurable \(Z\). Exhausting positive clearances proves
\eqref{eq:stopped-full-characteristic}, because \(R=h-M\) in the joint
limit. A countable determining family fixes one conditional version;
continuity then gives every real linear combination in the finite-vector
assertion, including independence across components that shared a
discrete precursor.

It remains to identify \(M\). On an active buffer in
\eqref{eq:stopped-partition}, \(M_\delta\) is exactly its flat chamber
label. The protected rational-label coordinates and their agreement along
compact connecting paths show that
\(M_\delta(\phi_\delta^\eta)\) converges to
\(\int_U m_O\phi^\eta\). Terms of zero limiting weight vanish by
tightness of the fixed \(M_\delta(\phi\psi_i)\) pairings. The joint extraction also identifies this
limit with \(M(\phi^\eta)\). We spell out the bound on the limiting
labels. A finite deterministic cover of \(Q\) by rational balls of
radius at most \(\rho/8\) can be chosen so that every point lies in one
of them. On \(d_\delta(Q)\ge\rho\), any initial chamber meeting \(Q\)
contains such a ball with clearance at least \(\rho/2\) from
\(F_{0,\delta}\). Thus the maximum of its initial labels is bounded by
the maximum of finitely many protected initial-label coordinates and
is tight. On this clearance event, a contributor bound \(N_\delta(Q)\le M\)
and an initial-label bound \(B\) bound every terminal chamber label meeting
\(Q\) by \(B+2aM+a\). Transfer this bound first to any finite family
of protected limiting buffers, using their joint convergence and a
slightly smaller clearance. The probability of its failure is at most
the initial count or label tail. Exhaust the countable rational buffers;
the same tail bound holds for the supremum of the limiting labels on
\(Q\setminus A\). Sending \(M,B\) to infinity proves that this supremum
is finite almost surely on each positive-clearance localization.

On these truncations the omitted finite contribution vanishes by
\eqref{eq:mean-neighborhood}; the omitted limiting integral vanishes
by the local label bound and \(|Q\cap F^{5\eta}|\to0\). Joint convergence
of the two \(M_\delta\) pairings transfers the finite omission bound
to \(M(\phi)-M(\phi^\eta)\) in probability. Letting \(\eta\downarrow0\)
therefore gives
\[
                 M(\phi)=\int_U m_O(z)\phi(z)\,\dd z
       \quad\hbox{on }\Lambda_Q.
\]
The initial whole-box tail of \cref{cor:global-cluster-counts} and label tightness remove the truncations
in probability. This also proves local \(\mathscr E\)-measurability of
\(M\). No mean or centered distribution remains on the cut in the
tested region. Countable rational localizations and a separate retained
state for each stop give the simultaneous assertion.
\end{proof}

\subsection{The stopped two-valued sets}

\begin{proposition}[Convergence of stopped carpets]
\label{prop:stopped-tvs}
The wired boundary cluster converges jointly with \(h_\delta\), with its
labels, to \(\TVS{a}{a}(h)\). In a free flat parent, after subtracting
its label and writing \(f\) for the zero-boundary limiting remainder,
the first odd carpet converges to \(\TVS{2a}{2a}(f)\). At depth \(r\),
define the global closed carpet as the closed complement of the holes
still unsearched after every branch has met \(r\) odd rims; it includes
all earlier layers and, in the wired case, the first wired layer. For each fixed
finite \(r\), this global carpet converges in ordinary Hausdorff distance
to the corresponding canonical nested carpet. Its labels and restarted
zero-boundary fields converge on every finite family of pointed compact
buffers in the limiting holes. These pointed conclusions include distinct
limiting components with the same discrete precursor; they do not assign
one whole precursor carpet to several separate Hausdorff limits.
\end{proposition}

\begin{proof}
Take a joint subsequence as in \eqref{eq:stopped-joint-extraction},
also retaining the raw closed carpets. Compactness of closed subsets,
field tightness, and the finite label sets permit this extraction.
The exact cuts give residual labels \(\pm a\) for the wired layer and
\(\pm2a\), relative to the parent, for the first free odd carpet.
Their closed limits are boundary-connected. Every open component is
represented by rational interior buffers: if a disk lay outside the
closed limit but in no pointed component, its compact containment in
the discrete complement would give a contradiction.

\Cref{lem:stopped-full-field} gives the full local-set decomposition
and its bounded label mean on tests throughout the parent. It also
gives zero area inside that parent. Hence the limit is a
boundary-connected bounded thin local set, including across its cut.
The uniqueness theorem of \cite[Proposition~2]{ASW2017} identifies
the two limits as \(\TVS{a}{a}\) and \(\TVS{2a}{2a}\), since
\(2a>2\lambda\) and \(4a>2\lambda\).

The original wired domain is Jordan. Its limiting holes are Jordan by the
two-valued-set geometry recalled above. The first odd holes have
label magnitudes \(2a>2\lambda\), so their loops avoid even a rough
parent boundary and are Jordan as well. Restart in such a hole after
subtracting its observed label. The conditional characteristic-function
estimate is uniform on deterministic interior buffers; a countable
exhaustion handles the random hole. Induction gives every fixed finite
odd depth, with a single global closed carpet at that depth and pointed
compact restrictions inside its holes. In the wired case, the local
applications are on compact subsets of pointed wired holes, and the
already identified \(W_0\) is supplied by the preceding layer.
Write \(K_r\) for the connected global limit at depth \(r\). If \(V\)
is a pointed Jordan hole of \(K_{r-1}\), then
\(\partial V\subset K_{r-1}\subset K_r\), and
\(K_r\cap\overline V\) is connected. Indeed, a separation of this
compact intersection would put the connected \(\partial V\) on one
side and leave the other side compactly inside \(V\). Adjoining
\(K_r\setminus V\) to the first side would then separate \(K_r\).
Thus the child restriction is boundary-connected. Inside each pointed
parent, uniqueness identifies that restriction; outside the union of
those parents the global closed set is the preceding layer. This identifies
the whole global closed limit even when one precursor splits. The sets
are canonical functions of the same limiting field, so every subsequential
law is the stated one.
\end{proof}

\subsection{Fields at seed stops}

\begin{proposition}[Conditional fields at a cluster stop]
\label{prop:exploration-thin}
Consider a seed stop in an initial free parent of label \(b_0\), or
the composite wired stop defined above. Let \(A,S,W_0,\mathscr E\)
be its limiting explored set, seed trace, initial wall, and retained
stopped sigma-field. The set \(A\) is connected to the original
boundary, and every component \(O\) of \(D\setminus A\) is simply
connected. Conditionally on \(\mathscr E\),
\begin{equation}
 h|_O=m_O+h^O,\qquad m_O-b_0(O)\in2a\Z,
 \label{eq:stopped-component-law}
\end{equation}
where the \(h^O\) are independent zero-boundary GFFs. Here \(b_0(O)\)
is the label of its initial parent chamber; it is \(0\) in the
initial free domain and \(\pm a\) in a wired hole. The labels are
tight on every finite family of pointed components.

On every compact subset of \(D\) away from \(W_0\cup S\), \(A\) has
zero area and the full conditional mean \(M\) is the locally
integrable piecewise constant function \(m_O\). On the same region,
\cref{eq:stopped-full-characteristic} gives the full zero-extended
conditional GFF law of \(h-M\), also for tests whose supports meet
\(A\). Locally \(M\) is the limit of the finite signed odd-hole
contributions of the discovered clusters, with the initial parent
label added. These statements hold for every member of a specified
countable family of sweeps and observable indexed stops, each
conditional only on its own retained state.

For a composite wired stop, every \(O\) is contained in one pointed
limiting wired hole. The conclusions use the global exploration of
the original \(h\); they do not require a separate lattice sweep
selected from that limiting hole.
\end{proposition}

\begin{proof}
The initial wall, the attached seed, and all clusters met by it form
one boundary-connected revealed union. Each residual face has one
sealed monochromatic \(W\) perimeter with its actual connections, by
\cref{lem:finite-cuts}. Its flat label differs from that of its initial
parent by an element of \(2a\Z\). Hausdorff limits preserve the
connected revealed union. The planar connected-complement criterion
then makes every open component simply connected. In the wired
composite case \(A\) contains \(W_0\), so each open component lies in
one component of \(D\setminus W_0\).

Before the limit, reflection of a centered flat law makes its mean
zero. Thus on an unsearched site the exact jump rule gives
\begin{equation}
 M_\delta(z)=b_{0,\delta}(z)+
       2a\sum_{C\ \mathrm{discovered}}
             \epsilon_C\ind_{\mathcal O_\delta(C)}(z),
 \qquad z\notin A_\delta,
 \label{eq:discrete-discovered-mean}
\end{equation}
where \(b_{0,\delta}(z)\) is its initial chamber label and
\(\epsilon_C\in\{-1,1\}\) is the cluster's relative jump sign.
The contributor bound and \cref{lem:cut-neighborhoods} give the
finite local truncation of this formula. On a connected compact
ball away from \(W_0\), the initial label is one constant. On exposed
sites the parity discrepancy from an incident wall is bounded and is
removed with the cut neighborhood.

\Cref{lem:stopped-full-field} now proves both the full mean identity
and the full residual law. Its interior restriction is
\eqref{eq:stopped-component-law}; the block Green kernel gives
independence even when one discrete chamber pinches into several
limiting components. The label tightness was proved in the joint
extraction. For the composite wired stop, all inputs are those of
the single global finite exploration, so the same argument covers
such a pinch without a future choice of lattice chamber.
\end{proof}

\subsection{Geometry of the candidate clusters}

The incoming label must be kept when a canonical cluster is restarted.
If its parent field is \(b+f\), the construction in
\cref{def:free-procedure} has the following values:
\[
\begin{array}{c|c|c|c}
\text{incoming label}&\text{CLE sign}&\text{field in the CLE hole}
     &\text{labels after the split}\\ \hline
b&+1&b+2\lambda+f_V&\{b,b+2a\}\\
b&-1&b-2\lambda+f_V&\{b-2a,b\}.
\end{array}
\]
Equivalently, \(b\) becomes \(b+\xi\,2\lambda\) and then one of
\(\{b,b+\xi\,2a\}\). A hole of a retained cluster is \emph{odd}
when its label is \(b+\xi\,2a\), and is even when its label is \(b\).
Thus every descendant baseline lies in \(b_0+2a\Z\). Here
\(b_0=0\) in the free domain and \(b_0=\pm a\) behind the wired layer.

We first justify the nested continuation that will be used twice.

\begin{lemma}[Completion between fixed labels]
\label{lem:nested-completion}
Let \(f\) be a zero-boundary GFF in a bounded simply connected parent
\(P\). Let \(E\) be a field-measurable boundary-connected bounded thin
local set whose value \(c_H\) is constant on each component \(H\) of
\(P\setminus E\). Suppose
\[
       -p\le c_H\le q,\qquad p,q>0,\qquad p+q\ge2\lambda.
\]
For these fixed \(p,q\), suppose also that all \(c_H\) lie in one fixed
deterministic countable subset \(\mathcal L\subset[-p,q]\).
In each component with \(-p<c_H<q\), insert
\(\TVS{p+c_H}{q-c_H}(f^H)\) in its zero-boundary remainder.
Leave components already labelled \(-p\) or \(q\) untouched.
The closed union is \(\TVS{p}{q}(f)\). In particular every such
partial continuation is contained in that same canonical target.
\end{lemma}

\begin{proof}
Point each component at its first rational point and use the normalized
Riemann map from the disk. The kernel theorem makes this map Borel in
the pointed domain. For each \(c\in\mathcal L\cap(-p,q)\), choose one
Borel version of the canonical \(\TVS{p+c}{q-c}\) map in the disk,
apply it to the pulled-back remainder, transport its interior set, and
adjoin the component boundary. This gives jointly measurable insertions
in the random pointed components for the countably many fixed pairs.
Enumerate the selected components by rational interior points and let
\(E_N\) insert only the first \(N\) complete two-valued sets.
Conditional local-set gluing in these distinct components makes
\(E_N\) local. Its mean distribution is the old mean plus the
zero-extended relative means of those insertions. On a new hole the
value is \(c_H\) plus its relative value, hence is \(-p\) or \(q\).
The old density on the inserted set contributes nothing, because a
two-valued set has zero area inside its open parent
\cite[Proposition~3.1]{AS2018}. This says nothing about the area of
the parent boundary. Thus \(E_N\) is thin, field-measurable,
boundary-connected, and a \(K\)-BTLS for \(K=\max(p,q)\).

Here is the needed justification for passing through countably many
holes. Choose an integer \(M\) with \(K\le2M\lambda\). By
\cite[Propositions~1 and~3]{ASW2017}, all \(E_N\) lie in the same
larger canonical set
\(Z=\TVS{2(M+1)\lambda}{2(M+1)\lambda}(f)\).
The compact-interior dimension bound of
\cite[Proposition~3.1]{AS2018} gives \(|Z\cap P|=0\) by exhaustion.
Let \(E_\infty=\overline{\bigcup_N E_N}\), with closure in
\(\overline P\). The increasing-local-set result
\cite[Lemma~2.3(3)]{ALS2017} applies: the sets are field-measurable,
and after adjoining \(\partial P\) each has one nontrivial connected
component. It makes \(E_\infty\) local and gives convergence of the
conditional mean distributions.

For every smooth compactly supported \(\varphi\), thinness gives
\(|f_{E_N}(\varphi)|\le K\|\varphi\|_{L^1(P)}\). Take an almost surely
convergent further subsequence on a countable dense family of tests,
pass this bound to the limiting distribution, and then use continuity.
By \(L^1\) duality that distribution has an
\(L^\infty(P)\) density bounded by \(K\). Locality identifies it
with the harmonic mean off \(E_\infty\). Since
\(E_\infty\cap P\subset Z\cap P\) has zero area, this is the full
mean distribution, so \(E_\infty\) is thin. Every complementary
component is an untouched terminal hole or a hole of one inserted
set: a compact subset of one original \(H\) cannot be affected by
insertions in the other components. Its value is therefore \(-p\)
or \(q\). Uniqueness in \cite[Proposition~2]{ASW2017} identifies
the union with \(\TVS{p}{q}(f)\).

This also proves the nested-interval rule for each fixed target interval
and a partial local set whose labels range over a fixed deterministic
countable set: after centering the incoming field, in \([L,U]\) with
\(L<0<U\) and \(U-L\ge2\lambda\), a hole of value \(c\in(L,U)\) is continued by the relative
set \(\TVS{c-L}{U-c}\), while a hole already at an endpoint is left
alone. The uses below have finite label sets. This is a same-field
statement, consistent with the monotonicity and nested construction in
\cite[Section~3.2]{AS2018}.
\end{proof}

The next lemma records the facts needed to identify how several odd
rims belong to one cluster. It does not assume local finiteness over
all generations.

\begin{lemma}[Canonical cluster geometry]
\label{lem:canonical-geometry}
In the free procedure of \cref{def:free-procedure} the retained sets
form a countable family of pairwise disjoint connected compact sets,
each strictly inside its parent domain. They and all their labelled
holes are measurable functions of the GFF. Each retained cluster is
the closure of the union of its odd-hole boundaries. Conditional on
the unlabelled recursive geometry, including which odd holes belong
to the same retained cluster, their relative jump signs are
independent and fair.

In a parent field \(b+f\), a positive first-generation retained
cluster is contained both in the positive excursion of \(f\) and in
the boundary-connected upper first-passage set of \(f\) at level
\(2a\). The containing cable sets have paths respectively above
\(b\) and below \(b+2a\) in total field units.
A negative retained cluster satisfies the reflected containments,
giving paths below \(b\) and above \(b-2a\). The carpet formed by
continuing through even holes and stopping at first odd holes is
\(\TVS{2a}{2a}(f)\). These statements remain valid at
\(a=2\lambda\).
\end{lemma}

\begin{proof}
Write \(P\) for the current parent. Each retained set is a two-valued set inside a Jordan $\CLE_4$
hole, with that hole's boundary included. Its closed set is
connected and compact. Distinct outer $\CLE_4$ loops are disjoint,
and every subsequent restart is strictly inside a complementary
hole of the previous retained set. Induction gives disjointness at
all finite generations; their union is countable because open
components each contain a rational point. Canonical measurability
follows at each restart from local-set measurability. None of these
observations asserts all-generation local finiteness.

Subtract the incoming \(b\) and condition on the outer CLE data.
Consider a positive \(\CLE_4\) hole \(V\). The field there is \(2\lambda+f_V\), and the
retained split is \(A=\TVS{2\lambda}{2a-2\lambda}(f_V)\), with
total values \(0\) and \(2a\). The first-passage set
\(F=A_{-2\lambda}(f_V)\) is the increasing limit obtained by sending
the upper two-valued barrier to infinity. Thus it contains \(A\),
and \(F\) together with the outer CLE rim is the positive excursion
inside that loop \cite{ALS2017,ALS2023}.

The FPS decomposition on \(V\) reads
\(2\lambda+f_V=\nu+f^F\), where
\(\nu=(f_V)_F+2\lambda\,\dd z\) is a positive measure whose relative
support is all of \(F\cap V\)
\cite[Definition~8 and Theorem~9]{ALS2023}.
The tower property for the nested sets \(A\subset F\) gives, for a
nonnegative smooth test \(q\) compactly supported in \(V\),
\begin{equation}
 \E[\nu(q)\mid V,A\text{ and their labels}]
       =2a\int_{\mathcal O(A)}q(z)\,\dd z,
 \label{eq:fps-odd-support}
\end{equation}
where \(\mathcal O(A)\) is the union of its odd holes. Indeed the
centered remainder after \(F\) has conditional mean zero, and thinness
of \(A\) makes the conditional mean of \(2\lambda+f_V\) exactly the
density \(2a\ind_{\mathcal O(A)}\). A countable family of nonnegative
tests, multiplied by the event that their supports miss the odd holes,
shows from positivity that \(\operatorname{supp}_V\nu\) is contained
in their closure. Since that support contains \(A\cap V\), every
interior neighborhood of a point of \(A\cap V\) meets an odd hole.
The point itself lies outside those interiors, so their boundaries
accumulate there.

The relative support statement does not include the outer rim.
Condition on \(V\) and fix a conformal map from the disk onto this
Jordan domain. On a closed angular subarc take a smooth
nonnegative probability density $w$ relative to $\dd\vartheta/(2\pi)$,
and push its averages on circles of radii $r\uparrow1$ into $V$.
Smooth radial approximations give compactly supported smooth tests
$q_r$ with the same limiting estimates. The disk Green kernel is
nonnegative, and its average against uniform angular measure on both
circles of radius $r$ is $-\log r$. Hence the Green energy of these
weighted averages is at most $\|w\|_\infty^2(-\log r)$, which tends
to zero. In particular $f_V(q_r)\to0$ in $L^2$.

The centered remainder \(f^F\) after the FPS continuation has
conditional covariance bounded above by $G_V$. Therefore
\(f^F(q_r)\to0\) in \(L^2\) as well. If an ambient neighborhood of
the chosen boundary arc missed the interior FPS, its positive measure
would vanish on $q_r$ for all sufficiently large $r$. On this event
the full FPS decomposition would give
\(2\lambda+f_V(q_r)=f^F(q_r)\).
Both centered terms tend to zero in probability, whereas their
difference on that event is the nonzero constant $-2\lambda$.
The event consequently has probability zero. A countable basis of
angular arcs and ambient neighborhoods, using the Jordan boundary
extension \cite[Chapter~I, Theorem~3.1]{GarnettMarshall2005}, proves
\(\partial V\subset\overline{F\cap V}\).
This is the boundary-average argument behind FPS nontriviality
\cite[Proposition~4.6]{ALS2017}, localized at the outer rim;
the closed excursion and its ambient support are also described in
\cite[Theorem~1 and Section~3]{ALS2023}.

By \cref{eq:fps-odd-support}, the interior FPS lies in the closure
of the odd holes. The preceding inclusion thus puts the outer rim
in that closure too. Since a point of $A$ is outside every odd-hole
interior, approaching it from an odd hole also approaches its
boundary. This proves the claimed closure identity on the whole
retained cluster.

For the upper FPS containment, expose the outer \(\CLE_4\) set and
this positive split. Finite conditional gluing makes it a
field-measurable boundary-connected BTLS with relative labels among
\(\{-2\lambda,0,2\lambda,2a\}\). Fix \(R>2\lambda\).
All but \(2a\) lie strictly between \(-R\) and \(2a\), since
\(a>\lambda\). By \cref{lem:nested-completion}, completing these
holes to \(-R,2a\) yields \(\TVS{R}{2a}(f)\), containing the split.
As \(R\uparrow\infty\) through a deterministic sequence, reflection of
the FPS construction
\cite[Remark~4.2 and Proposition~4.4]{ALS2017} identifies the closed
increasing union as the upper FPS of \(f\) at \(2a\). This last set
is generally nonthin; the preceding bounded-density argument is not
used as the barrier grows. Its metric approximants connect points to
the parent boundary by sublevel paths. Excursion and FPS convergence
joint with the field are supplied by \cite{ALS2018,ALS2023}.
The next section uses their connected paths between strict interior
hits. In the total field the two path levels are \(b\) and \(b+2a\).
Reflection gives levels \(b\) and \(b-2a\) for a negative cluster.

For the sign assertion, let $\xi$ be the fair outer $\CLE_4$ sign.
The zero-boundary remainder $f_V$ is independent of it. The field
$\xi f_V$ has the same law and is still independent of $\xi$.
Reflection carries the negative split into the positive split,
preserving its unlabelled geometry and the distinction between odd
and even holes. Conditional remainders in the holes can be
reflected in the same way before restarting. This proves the
independent-fair-sign assertion against every finite recursive geometry
cylinder. A monotone-class argument gives the assertion conditional on
the entire countable unlabelled geometry. It does not assert independent
labels on arbitrary two-valued-set loops.

Finally, fix the incoming \(b\), subtract it, and put
\(T=\TVS{2a}{2a}(f)\). At a completed first-odd stage, holes labelled
\(\pm2a\) are terminal and holes labelled zero are active. In every
active hole insert its complete outer CLE set, then the appropriate
complete split in each CLE hole. Define each layer by enumerating its
holes and taking the closed limit of finite insertions. Every finite
partial layer is a field-measurable boundary-connected BTLS by finite
gluing. Its labels lie among \(0,\pm2\lambda,\pm2a\), hence inside
\([-2a,2a]\). Its completion under \cref{lem:nested-completion} is
the same \(T\), so the partial layer lies in \(T\). For a full layer,
the increasing-local-set and bounded-density proof of that lemma applies
with \(Z=T\) and \(K=2a\). It remains a BTLS contained in \(T\).

Let \(B_0=\partial P\), let \(B_n\) be the completed stage after \(n\)
rounds of even continuation, and
let \(B_\infty=\overline{\bigcup_n B_n}\). The same increasing-local-set
result makes \(B_\infty\) local and gives convergence of the mean
distributions. Their absolute values on a test \(\varphi\) are bounded
by \(2a\|\varphi\|_1\), so the limiting mean has an \(L^\infty\)
density bounded by \(2a\). Since \(B_\infty\subset T\) and
\(|T\cap P|=0\), that density agrees with the harmonic mean as a full
distribution. This proves thinness of the increasing limit.

The active region has vanishing area. Conditional on a positive CLE
hole, thinness and mean zero for \(f_V\) give, for Lebesgue-a.e. \(z\),
\[
 0=-2\lambda(1-p(z))+(2a-2\lambda)p(z),
 \qquad p(z)=\lambda/a,
\]
where \(p(z)\) is the conditional probability that the split assigns
the odd value. Reflection gives the same probability in a negative
hole. Fubini therefore bounds the expected area still active after
\(n\) rounds by \((1-\lambda/a)^n|P|\). Monotone convergence makes
the area surviving all rounds zero almost surely.
Terminal labels persist, so the limiting harmonic mean takes values
\(\pm2a\) almost everywhere off \(B_\infty\). Harmonic continuity and
connectedness make it one of these constants throughout each open
component. BTLS uniqueness gives \(B_\infty=T\). At the next odd
depth the observed \(b\pm2a\) becomes the new incoming label before
the argument is restarted.

At \(a=2\lambda\) the split has width \(4\lambda\) and both relative
barriers remain positive. The FPS support proof
\eqref{eq:fps-odd-support}, the sign argument, and the containments
still apply, and the odd probability is \(1/2\). Thus the endpoint is
included.
\end{proof}
